\documentclass[openany,fontset=fandol]{ctexbook}
\usepackage{amsmath,amssymb}
\usepackage[unicode,hidelinks]{hyperref}
\setlength{\emergencystretch}{2em}
\interfootnotelinepenalty=10000
\raggedbottom
\DeclareMathOperator{\Div}{div}
\DeclareMathOperator{\Ens}{Ens}
\newcommand{\OO}{\mathcal O}
\newcommand{\KK}{\mathcal K}
\newcommand{\PP}{\mathbf P}
\begin{document}
\title{代数几何基础第五卷预备笔记\\线性系的除子集合与底除子}
\author{据法文扫描和英译编辑稿互校续译}
\date{2026年11月22日}
\maketitle
\chapter*{范围与接续}
第85批续译第三篇命题4.4及其后作者说明，至第三篇末。
法文扫描物理146--147页（手写153--154），英译编辑稿物理110--111页。
命题的 a) 保留原意；b) 的字面断言有反例，c) 也随之补足平移条件。正文给出经补证的修订表述，
脚注保留手稿原意与反例。作者末段按原稿译出，其讨论范围在脚注注明。
下一批从法文148页、英译112页第四篇致迪厄多内的信起。
本批未并入1158页合订本，不改动正式出版的 EGA IV.1。

\chapter*{第三篇（续完）}
\section*{4\quad 除子线性系与可逆模（续）}
再回到地面上，甚至再贴近地面一些：为了用通俗的话说明线性系的概念，
姑且取一个基域以求简便（尽管换成仿射底，陈述实质上仍可保持原样）。\footnote{法146的基域为 \emph{corps de base}；英110的“box field”非数学术语。关于仿射底的话是作者的提示，不是此处已经证出的推广。}

\medskip\noindent\textbf{命题4.4。}
设 \(X\) 是域 \(k\) 上有限型的预概形，且自然映射
\[
 k\longrightarrow H^0(X,\OO_X)
\]
是同构。对 \(X\) 上由 \(k\) 上射影丛 \(P^\vee\) 参数化的每个除子线性系 \(D\)，
记 \(\Ens(D)\) 为形如 \(D(t)\) 的除子组成的集合（注意：是集合！），其中
\(t\in P^\vee(k)\)。\footnote{\(\Ens(D)\) 只取 \(k\)-有理参数，归入 \(X\) 上的 Cartier 除子集合；不是取线性等价类，也不是命题4.3中的除子子函子。若允许零维参数向量空间，约定其射影空间为空；若只定义正秩射影丛，则以下线性系的除子集合非空。}

\smallskip\noindent\textbf{a)}
设 \(D'\) 是 \(X\) 上另一个除子线性系，由射影丛 \((P')^\vee\) 参数化。
则 \(D\) 与 \(D'\) 同构，当且仅当 \(\Ens(D)=\Ens(D')\)。

\smallskip\noindent\textbf{b)（校订）。}
假设 \(k\) 代数闭，或 \(X\) 几何整。\footnote{法146边注补入此二选一条件；b) 随后的“为空，或……”被划去，并改接法147的向量子空间条件。此处不将删文复作正文。英111保留了边注的两项假设，但省略了若干手稿改写痕迹。}
记 \(\KK_X\) 为全分式层。\(X\) 上非空的有效 Cartier 除子集合 \(\Delta\) 形如
\(\Ens(D)\)，当且仅当存在有效 Cartier 除子 \(A\) 和有限维非零
\(k\)-向量子空间 \(E\subseteq\Gamma(X,\KK_X)\)，使每个 \(f\in E\setminus\{0\}\)
都属于 \(\Gamma(X,\KK_X^\times)\)，且
\[
 \Delta=\{\, A+\Div(f): f\in E\setminus\{0\}\,\}.
\]
右侧中的每个除子必须有效。可以取 \(A\in\Delta\)；以它为底时，可以取 \(1\in E\)。
这里 \(A\) 只是用来比较的基准除子，不必是线性系共有的固定部分；\(\Div(f)\) 是主 Cartier 除子，允许有极点。\footnote{法147改称亚纯函数，\emph{régulière} 由“\(\Div(f)\) 有定义”解释，即在全分式层中可逆；不是说函数处处无极点。手稿与英111都把 \(\Delta\) 直接写成 \(\{\Div(f)\}\)，未写平移项，也未明确限制 \(E\) 有限维。按字面，若 \(\Div(f)\) 有效，\(f\) 局部是非零因子与单位的乘积，故 \(f\in H^0(X,\OO_X)=k\)，只能给零除子。反例：\(\PP^1_k\) 上 \(\OO(1)\) 的完全线性系包含一次点除子。正文采用可由4.1严格推出的修订，不把它冒称手稿原句。}

\smallskip\noindent\textbf{c)（校订）。}
设 \((A,E)\)、\((A',E')\) 满足 b) 的条件，并分别给出 \(\Delta\)、\(\Delta'\)。
则 \(\Delta=\Delta'\)，当且仅当存在 \(h\in\Gamma(X,\KK_X^\times)\)，使
\[
 E'=hE,\qquad A'=A-\Div(h).
\]
存在时，\(h\) 在乘以 \(k^\times\) 的元素的意义下唯一。\footnote{法147的 c) 仍印有“严格有理函数”，对乘子又改作“正则亚纯函数”；英111相应使用 \emph{regular pseudo-function}，并误排 \(k^\times\)。手稿只写乘法条件 \(E'=hE\)，未写除子的抵消关系。若严格要求两组主除子都有效，则 b) 的字面表述只剩平凡情形。若放宽为一般主除子来理解 c) 的乘法断言，取 \(X=\PP^1_k\)、\(E=k\)、\(E'=kx\)、\(h=x\)，则两个集合分别是 \(\{0\}\) 与 \(\{[0]-[\infty]\}\)。反过来，不带平移的主除子集合若相等，则每对同除子函数之比是全局单位，在本命题下属 \(k^\times\)，所以 \(E=E'\)。}

\medskip\noindent\textbf{作者说明（译）。}
利用4.1，证明就是一个简单的练习；除非你要求，我就不写出来了。
此外我觉得，把4.4移到4.3前面也许更好，因为在技术上它要简单得多。
还应注意，若 \(X\) 几何整，b) 对 \(E\) 所提的那项条件自动满足。\footnote{指非零亚纯函数为全分式层的单位这一条件；对于校订后的表述，
有限维性以及 \(A+\Div(f)\) 的有效性仍须满足。}

在 b) 开头加限制，是因为否则，b) 所述条件在过到 \(k\) 的代数闭包后可能不再成立。
在任一特征中都可以举例，甚至 \(k\) 在正特征下可分闭时也这样。
要把话说周全，b) 应不附加关于 \(k\) 或 \(X\) 的条件，而是在过到 \(k\) 的代数闭包后
陈述对 \(E\) 的条件（并指出在几何整情形此项条件自动满足）。
稍微滥用一下术语，形如 \(\Ens(D)\) 的除子集合 \(\Delta\)，常径称为 \(X\) 上的除子线性系。\footnote{最后一句为法147手写；英111标作略难辨读，并在其后加了手稿正文告终的编辑说明和拉丁文结束语。法148另起书信。这里的“在代数闭包上检验”是检验基变换后的向量空间及其全部新线性组合，不是只检查原有 \(k\)-有理参数。}

\clearpage
\section*{校订补证（非原文）}
\noindent\textbf{1. 从除子到截面。}
有效 Cartier 除子 \(B\) 带有可逆层 \(\OO_X(B)\) 及其典范正则截面 \(s_B\)。
反之，可逆层 \(L\) 的正则截面 \(s\) 给出零除子 \(B\)，以及唯一把 \(s_B\)
送到 \(s\) 的同构 \(\OO_X(B)\simeq L\)。局部若 \(B\) 由非零因子 \(b\)
定义，则 \(\OO_X(B)\) 在全分式层中由 \(1/b\) 生成，\(s_B\) 对应函数 \(1\)。
因此亚纯函数 \(f\) 若是全分式层的单位，则 \(f\) 给出 \(\OO_X(B)\)
的全局正则截面，当且仅当 \(B+\Div(f)\) 有效：局部等价于 \(bf\)
是正则函数且为非零因子。\footnote{现代记法及正则截面与有效 Cartier 除子的对应见
\href{https://stacks.math.columbia.edu/tag/0C4S}{Stacks Project，0C4S}。
这里“截面正则”指乘法给出单射，不意味着截面处处不消失。}

\medskip\noindent\textbf{2. 证明 a)。}
由4.1，正秩线性系可由 \((L,V)\) 表示，其中 \(V\subseteq H^0(X,L)\)
有限维，且在几何基变换后，\(V\) 的每个非零向量仍给正则截面。参数是 \(V\)
中一维子空间；这是原稿的对偶射影丛记法。若两个除子集合相同，取一个公共除子
\(B\)，选取产生它的截面 \(s\in V\)、\(s'\in V'\)，并用第一步把两个可逆层都
识别为 \(\OO_X(B)\)，使所选截面都对应 \(s_B\)。在同一可逆层上，产生同一
除子的两个正则截面局部相差单位。这些单位唯一，故粘为全局单位，在本题中
属 \(k^\times\)。于是任取 \(0\ne v\in V\)，可从 \(V'\) 选出产生同一除子的
截面，该截面是 \(v\) 的非零标量倍数，故 \(v\in V'\)；交换二者得 \(V=V'\)。
由4.1恢复的参数空间及通用除子相同，故线性系同构。反向由同构对
\(k\)-值点给出双射。若容许空参数，\(k\) 上任意正秩射影空间均有 \(k\)-点，
故两个空集合只对应同一个空参数情形。

\medskip\noindent\textbf{3. 证明 b)。}
给定线性系，选一非零参数截面 \(s_A\)，以其除子为 \(A\)，并把可逆层识别为
\(\OO_X(A)\)。每个非零截面除以 \(s_A\) 得到全分式层的单位；它们与零合成一个
有限维向量子空间 \(E\)，包含 \(1\)。截面的除子正是 \(A+\Div(f)\)，所以
得到右侧的表示。

反过来，由第一步，\(E\) 是 \(H^0(X,\OO_X(A))\) 的有限维子空间，且其中
非零元素给正则截面。若 \(k\) 代数闭，这已是在4.1所需的代数闭包上的检验。
若 \(X\) 几何整，则 \(X_{\overline k}\) 为整概形，扩域后的任何非零向量仍给
非零截面，因而正则。这里用到
\[
 H^0(X,L)\otimes_k\overline k\simeq H^0(X_{\overline k},L_{\overline k}).
\]
有限型的域上预概形拟紧拟分离，取有限仿射覆盖，并用有限仿射开覆盖重叠部分，
可把全局截面写成两条有限乘积间映射的等化子；平坦扩域与这一等化子交换。
于是4.1给出所需的线性系。若无 b) 开头的二选一条件，同一证明仍成立，
但应直接要求几何基变换后所有非零向量给正则截面。

\medskip\noindent\textbf{4. 证明 c)。}
若 \(\Delta=\Delta'\)，取公共除子 \(B=A+\Div(f_0)=A'+\Div(f'_0)\)。
把两套数据都改用 \(B\) 表示，所用函数空间分别为 \(E/f_0\) 与 \(E'/f'_0\)。
它们是在同一可逆层 \(\OO_X(B)\) 中产生同一除子集合的两个截面空间，
故由 a) 的逐截面比较相等。取 \(h=f'_0/f_0\) 即得
\(E'=hE\) 且 \(A'=A-\Div(h)\)。反过来这两式使每个
\(A'+\Div(hf)\) 恰等于 \(A+\Div(f)\)，故集合相同。
若 \(h_1,h_2\) 都给出这两式，\(\Div(h_2/h_1)=0\)，所以
\(h_2/h_1\) 是全局单位，即属 \(k^\times\)；乘任何非零常数也不改变这两式。
\end{document}
